\section{Discussion}\label{sec:discussion}
This chapter discusses the findings and the limitations of the study. The study evaluated whether the price paths generated by a GAN reproduce the implied volatility of the market. They do not. The generated prices deviate further from the market than the prices of GBM and of both bootstrap modes in every moneyness range. In addition, the deviation is a systematic underpricing. The result holds for ForGAN on TTF front-month futures over the test period of Sec.~\ref{sec:data-split}.

\subsection{Sources of the Deviation}\label{sec:deviation-sources}

Several factors contribute to the deviation between the generated option prices and the market prices. These include the spread of the generated distribution, the size of the training sample, the returns used as the condition and the recursive rollout of Sec.~\ref{sec:pricing-setup}.

The generated distribution is narrower than the returns of the test split. The standard deviation of the generated returns is \num{0.0369} against \num{0.0404} for the test split, and the excess kurtosis \num{10.12} against \num{13.27} (Sec.~\ref{sec:network-capacity}). As a result, the option prices are underestimated across all moneyness ranges. The narrower distribution is not caused by mode collapse, since no
one-step distribution of $G$ collapsed (Sec.~\ref{sec:baseline}).
The underpricing is visible not only in the tails but also at the at-the-money strikes. $G$ has an implied volatility error of \num{0.243} for at-the-money strikes, against \num{0.144} for the block bootstrap. The prices at the at-the-money strikes do not depend on the extreme returns, so the deviation is not caused by the tail weight alone. 

The size of the training sample limits the tail weight of the generated returns. The training split contains \num{3197} returns (Sec.~\ref{sec:data-split}). \textcite{fuTimeSeriesSimulation2019} report that a cGAN underestimated the excess kurtosis of a simulated series at a training sample of \num{1000} observations and matched it at \num{20000}. The excess kurtosis of the generated returns is below the test split, which is consistent with that finding. The available TTF futures data covers the period from 2010 (Sec.~\ref{sec:futures-data}), so a longer training sample would require additional data.

The condition variable used in the GAN is the $l$ past returns, so $G$ has to infer the volatility from them. As seen in Fig.~\ref{fig:conditioning}, $G$ responds to the conditioning. The spread of the generated returns correlates with the spread of the condition window at \num{0.752}. In contrast, the spread of the bootstrap stays almost constant across the \num{20} valuation dates. However, the spread of $G$ stays below the implied volatility of the market on \num{17} of the \num{20} valuation dates. \textcite{fuTimeSeriesSimulation2019} propose adding the LSTM cell to the generator so that the historical volatility is extracted from the condition, which ForGAN already contains (Sec.~\ref{sec:forgan}). The response of $G$ to the condition window shows that the volatility is extracted, but this does not translate into a match of the implied volatility. A longer condition window does not close the gap. On the contrary, increasing $l$ from \num{10} to \num{30} reduces the tail weight of the generated returns further (Sec.~\ref{sec:condition-window}). 

The generated returns are produced recursively (Sec.~\ref{sec:pricing-setup}). After $l$ steps, the condition window holds only the generated returns. At this point, $G$ operates outside the regime it was trained on. At moneyness 1.5 to 2.0, the share of inverted strikes is \SI{53.8}{\percent}, with a seed spread of \num{43.3} percentage points (Tab.~\ref{tab:price-error-generator}). The seeds disagree about whether the paths reach the outer strikes at all. The seed spread reported in Sec.~\ref{sec:network-capacity} is measured one step ahead and does not reflect the spread after $h$ recursive steps. The rollout was evaluated only at the end of the path, so the contribution of the recursive rollout to the deviation is not established.

\subsection{Limitations}\label{sec:limitations}

The results are subject to four limitations: the model selection criterion, the coverage of the benchmark set, the exclusion of the strikes that carry no implied volatility and the composition of the test split.

The model was selected on the test split, which is also used for the option pricing evaluation. The ForGAN baseline model was chosen on the seed spread of the generated returns, following the selection of configuration 8 in Sec.~\ref{sec:network-capacity}. The $\mathrm{PnL}^{*}_{a}$ and $\mathrm{PnL}^{*}_{a}$ \& $\mathrm{MSE}$ branches have a narrower seed spread than the ForGAN baseline on most of the distributional statistics of Sec.~\ref{sec:effect-cost}, with the excess kurtosis as the exception. However, their seed spread is wider on the implied volatility error of Sec.~\ref{sec:option-prices}. As a result, the selection criterion does not predict the seed spread of the option prices.

The benchmark coverage is limited to GBM and the two bootstrap modes. GARCH with $t$-distribution innovations \parencite{bollerslevConditionallyHeteroskedasticTime1987} was not included, since fitting it requires choices of the model order and the innovation distribution, which the bootstrap avoids. GARCH-$t$ conditions the simulated variance $\sigma^2_t$ on the returns preceding the valuation date, which neither bootstrap mode nor GBM does. The study does not establish whether the generated prices improve on a conditioned benchmark. GARCH-$t$ also reproduces the heavy tails without the asymmetry, which no other benchmark in the set does (Sec.~\ref{sec:benchmarks}).  The effect of asymmetry and heavy tails on the option prices is not separated.

The strikes with no implied volatility are excluded from the errors of Sec.~\ref{sec:option-prices}. These are the strikes that lie outside the range the simulated paths cover (Sec.~\ref{sec:pricing-setup}). As a result, the reported errors are lower bounds for $G$ and for GBM. The moneyness band from \num{0.5} to \num{2.0} (Sec.~\ref{sec:options-data}) includes strikes that the simulated distribution does not reach. A narrower band would raise the percentage of inverted strikes but would also discard the strikes where the models differ most.

The test split covers \num{401} trading days. The 2026 disruption (Sec.~\ref{sec:gas-markets}) includes \num{57} of these days. Excluding the disruption window reduces the excess kurtosis from \num{13.24} to \num{1.57} and the skewness from \num{0.91} to \num{-0.06}. As a result, the distributional comparison of Chapter~\ref{sec:exp} and the pricing comparison of Sec.~\ref{sec:option-prices} depend on a single event. Among the valuation dates of Sec.~\ref{sec:pricing-setup}, 2026-03-25 falls within the disruption. On this date, the market implied volatility reflects the expected price moves, while both bootstrap modes and $G$ are based on past returns (Fig.~\ref{fig:vol-smile-models}).